% Preliminary follow-up evaluation for Intent2Proof proposal.
% Include this only as SOURCE-GROUNDED evidence until real-RTL SBY tasks finish.
\subsection{Abstraction qualification: heterogeneous AXI-Lite control}
A follow-up to our open-source AXI-Lite benchmark examines whether the smaller
contract candidate preserves the observable accepted-request events of the
pinned upstream interconnect. We found that the concrete implementation uses
an initialized register for its read-address-ready output, whereas the candidate
uses a combinational ready expression when idle. Under source-grounded
post-reset conditions, Z3 returns \texttt{SAT} for a mismatch with either one
read requester or two simultaneous requesters, and \texttt{UNSAT} for a
no-request mismatch. These are one-step checks of source-anchored equations,
\emph{not} executed equivalence checks on fully elaborated RTL.

This result exposes a limitation of the earlier benchmark: two models may
independently satisfy the same safety property even if their handshake traces
differ. We therefore prepared a matched-input SymbiYosys miter that compares
accepted read-address events and a separate coverage property requiring both
masters to present requests to the same slave and eventually complete their
transactions. Negative controls introduce response-routing and unsafe-shutdown
faults. The eight prepared real-RTL tasks have not yet been run in this
execution environment. The run gate will reject the current candidate as a
cycle-exact replacement if the expected miter witness is confirmed and will
\emph{not} treat a shorter verification run as evidence of sound abstraction.

The strongest product implication is the need to \emph{qualify} every
property-directed abstraction against a precise observation alphabet and
assumption set before reusing its proofs for the actual FPGA build.
